%%%%%%%%%%%%%%%%%%%%%%%%%%%%%%%%%%%%%%%%%%%%%%%%%%%%%%%%%%%%%
%  Armaan Sachdeva — Master CV
%  Compiled with: pdflatex armaan_sachdeva_cv.tex
%%%%%%%%%%%%%%%%%%%%%%%%%%%%%%%%%%%%%%%%%%%%%%%%%%%%%%%%%%%%%
\documentclass[10pt, a4paper]{article}

% ── Packages ────────────────────────────────────────────────
\usepackage[a4paper, top=1.8cm, bottom=1.8cm, left=1.8cm, right=1.8cm]{geometry}
\usepackage{parskip}
\usepackage{titlesec}
\usepackage{enumitem}
\usepackage{hyperref}
\usepackage{xcolor}
\usepackage{fontawesome5}
\usepackage{array}
\usepackage{tabularx}
\usepackage{booktabs}
\usepackage{microtype}
\usepackage{multicol}
\usepackage{amsmath}
\usepackage{amsfonts}
\usepackage{lmodern}
\usepackage[T1]{fontenc}
\usepackage[utf8]{inputenc}

% ── Colours ─────────────────────────────────────────────────
\definecolor{accent}{HTML}{1A3A5C}      % deep navy
\definecolor{rule}{HTML}{2E6DA4}        % mid blue for rules
\definecolor{light}{HTML}{5A7FA8}       % lighter blue for sub-text
\definecolor{linkcolor}{HTML}{1A6496}

% ── Hyperref ────────────────────────────────────────────────
\hypersetup{
  colorlinks=true,
  urlcolor=linkcolor,
  linkcolor=linkcolor,
  pdftitle={Armaan Sachdeva — CV},
  pdfauthor={Armaan Sachdeva}
}

% ── Section formatting ──────────────────────────────────────
\titleformat{\section}
  {\large\bfseries\color{accent}}
  {}{0em}
  {}
  [\color{rule}\titlerule]

\titlespacing{\section}{0pt}{10pt}{4pt}

% ── List settings ───────────────────────────────────────────
\setlist[itemize]{leftmargin=1.4em, topsep=2pt, itemsep=1pt, parsep=0pt}
\setlist[itemize,1]{label=\textcolor{rule}{\textbullet}}

% ── Custom commands ─────────────────────────────────────────
\newcommand{\cvname}[1]{%
  {\Huge\bfseries\color{accent} #1}%
}
\newcommand{\cvtitle}[1]{%
  {\large\color{light} #1}%
}

% project entry: {Title}{Tech tags}{bullets}
\newenvironment{project}[2]{%
  \textbf{\color{accent}#1} \hfill
  \small\textit{\color{light}#2}%
  \begin{itemize}%
}{%
  \end{itemize}%
  \vspace{2pt}%
}

% ── No paragraph indent ─────────────────────────────────────
\setlength{\parindent}{0pt}

%%%%%%%%%%%%%%%%%%%%%%%%%%%%%%%%%%%%%%%%%%%%%%%%%%%%%%%%%%%%%
\begin{document}
\pagestyle{plain}

% ════════════════════════════════════════════════════════════
%  HEADER
% ════════════════════════════════════════════════════════════
\begin{center}
  \cvname{Armaan Sachdeva}\\[4pt]
  \cvtitle{3rd-Year Physics Undergraduate}\\[6pt]
  \small
  \faEnvelope\;\href{mailto:as45g22@soton.ac.uk}{as45g22@soton.ac.uk}\quad
  \faGithub\;\href{https://github.com/BBCnewslondon}{github.com/BBCnewslondon}\quad
  \faLinkedin\;\href{https://www.linkedin.com/in/armaan-sachdeva-79b68b214/}{linkedin.com/in/armaan-sachdeva-79b68b214}\quad
  \faMapMarker*\;United Kingdom
\end{center}

\vspace{2pt}
{\color{rule}\hrule height 0.8pt}
\vspace{6pt}

% ════════════════════════════════════════════════════════════
%  PROFILE
% ════════════════════════════════════════════════════════════
\section{Profile}

High-achieving Physics undergraduate transitioning into a Master's degree, with a demonstrable record of
translating rigorous mathematical and physical theory into production-quality software.
Expertise spans \textbf{Computational Astrophysics}, \textbf{Machine Learning}, \textbf{Numerical Simulation},
\textbf{Full-Stack Engineering}, and \textbf{Quantitative Finance}.
Seeking a summer internship where deep physical intuition and software engineering proficiency can be applied
to real-world computational challenges.

% ════════════════════════════════════════════════════════════
%  EDUCATION
% ════════════════════════════════════════════════════════════
\section{Education}

\textbf{BSc (Hons) Physics} \hfill \textit{2022\,--\,2025}\\
\textit{University of Southampton} \hfill \textcolor{light}{Expected 2:1}

\begin{itemize}
  \item Relevant modules: Quantum Mechanics, General Relativity, Electrodynamics, Statistical Mechanics,
        Computational Physics, Advanced Mathematical Methods, Nuclear \& Particle Physics.
\end{itemize}

\vspace{4pt}
\textbf{A-Levels} \hfill \textit{2020\,--\,2022}\\
Physics (A), Mathematics (A*), Geography (B)

% ════════════════════════════════════════════════════════════
%  DISSERTATION
% ════════════════════════════════════════════════════════════
\section{Dissertation --- \textit{``Ripples in Reality: The Theoretical Origin and\\
Observational Confirmation of Gravitational Waves''}}

\begin{itemize}
  \item \textbf{Mathematically Derived} the linearised perturbation of the Einstein field equations
        $G_{\mu\nu}=8\pi T_{\mu\nu}$ to obtain the gravitational-wave equation
        $\Box\bar{h}_{\mu\nu}=-16\pi T_{\mu\nu}$ in the Lorenz gauge, establishing the wave solutions
        $h_{\mu\nu}^{\mathrm{TT}}$ in the transverse-traceless frame.
  \item \textbf{Implemented} the quadrupole formula
        $h_{ij}^{\mathrm{TT}}=\tfrac{2G}{rc^4}\ddot{I}_{ij}^{\mathrm{TT}}$
        numerically to predict the chirp waveform for a compact binary inspiral,
        matching the morphology of the GW150914 LIGO detection with a
        template-fitting residual $<5\%$.
  \item \textbf{Visualised} time-frequency spectrograms of the GW150914 strain data using SciPy
        and Matplotlib, overlaying matched-filter templates generated from numerical relativity
        waveform catalogues (LALSuite-compatible), and provided a critical analysis of detector
        noise budgets, calibration uncertainties, and the statistical significance ($5.1\sigma$) of
        the detection.
\end{itemize}

% ════════════════════════════════════════════════════════════
%  TECHNICAL SKILLS
% ════════════════════════════════════════════════════════════
\section{Technical Skills}

\begin{description}[leftmargin=4.4cm, style=sameline, labelsep=0.4em,
                    font=\bfseries\color{accent}, itemsep=3pt, parsep=0pt, topsep=2pt]
  \item[Programming Languages]
    Python (NumPy, SciPy, Matplotlib, pandas, PyTorch, TensorFlow, Optuna),
    C++17 (CMake, STL), Rust, TypeScript/JavaScript (Node.js, React, Next.js),
    Dart (Flutter), SQL, C, HTML/CSS
  \item[Physics \& Maths]
    Numerical Integration (RK4, Dormand-Prince DOP853), Finite Difference Methods
    (explicit \& implicit), General Relativity \& Linearised Gravity, Quantum Mechanics,
    PDEs, Monte Carlo Simulations, Lattice Boltzmann Methods,
    Variational Methods, Statistical Mechanics
  \item[Machine Learning \& AI]
    Physics-Informed Neural Networks (PINNs), Convolutional Autoencoders (1D CNN, U-Net),
    Recurrent Architectures (LSTM, WaveNet), Hidden Markov Models (HMM),
    Hyperparameter Optimisation (Optuna), OpenCV, Gemini API
  \item[Web \& Mobile]
    React Native, Flutter, Next.js, Vite, Tailwind CSS, Socket.io, Flask, Capacitor
  \item[Databases \& DevOps]
    PostgreSQL, PostGIS, SQLite, Docker, Git, CMake, Linux (Bash)
  \item[Scientific Libraries]
    Uproot, Awkward Array, ROOT (via Uproot), LALSuite-compatible waveforms,
    face\_recognition, Streamlit
  \item[Domain APIs]
    OANDA REST API, Google Maps Platform, OpenRouteService, OR-Tools, Google Gemini 2.0
  \item[Authoring]
    \LaTeX, Markdown, technical report writing
\end{description}

% ════════════════════════════════════════════════════════════
%  PROJECTS — COMPUTATIONAL ASTROPHYSICS & RELATIVITY
% ════════════════════════════════════════════════════════════
\section{Projects: Computational Astrophysics \& Relativity}

\begin{project}{Black-Hole Ray Tracer \& Gravitational Lensing Visualiser}{C++17, Physically-Based Rendering, General Relativity}
  \item \textbf{Engineered} a C++ ray tracer that renders a black hole scene featuring a relativistic accretion disk, event horizon, and gravitational lensing.
  \item \textbf{Implemented} relativistic effects including Doppler beaming and light-bending, with animation capabilities to simulate a 120-frame orbital sequence at a safe distance.
  \item \textbf{Optimised} rendering with multi-threading and a custom vector math library (\texttt{vec3.h}), supporting PPM image output and FFmpeg-based video generation.
\end{project}

\begin{project}{Gravitational Wave Chirp Visualisation}{Python, NumPy, Matplotlib, GWPy, GWOSC}
  \item \textbf{Simulated} a Newtonian inspiral chirp (GW150914-like) to generate simulated strain and instantaneous frequency plots.
  \item \textbf{Implemented} dynamic amplitude scaling to overlay simulated waveforms on whitened public LIGO data for visual comparison.
  \item \textbf{Visualised} data using Python (NumPy, Matplotlib, and GWPy) to fetch and whiten real GWOSC signals.
\end{project}

\begin{project}{N-Body Gravitational Simulator (DOP853 \& Adaptive Stepping)}{Python, NumPy, SciPy, Matplotlib}
  \item \textbf{Developed} a comprehensive simulation suite for two-body and three-body problems using the high-order Dormand-Prince (DOP853) numerical integrator.
  \item \textbf{Implemented} adaptive step-size control and collision detection to maintain machine-precision energy conservation (relative error $\sim 10^{-15}$).
  \item \textbf{Visualised} chaotic dynamics including the Figure-8 orbit and hierarchical systems via 6-panel analysis plots, phase space diagrams, and real-time animations.
\end{project}

% ════════════════════════════════════════════════════════════
%  PROJECTS — MACHINE LEARNING & AI
% ════════════════════════════════════════════════════════════
\section{Projects: Machine Learning \& AI in Physics}

\begin{project}{Astrophysical Signal Denoising --- ``The Clarifyer''}{Python, PyTorch, TensorFlow, 1D CNN, U-Net, LSTM, WaveNet}
  \item \textbf{Trained} neural network autoencoders (1D CNN, U-Net, LSTM, and WaveNet-inspired) to denoise simulated gravitational-wave and pulsar signals.
  \item \textbf{Engineered} a noise modelling suite incorporating Gaussian, coloured detector noise (LIGO PSD), and instrumental glitches (blips, whistles).
  \item \textbf{Optimised} performance using PyTorch and TensorFlow, achieving SNR improvements of 5--15\,dB and correlation coefficients $>0.8$ with original signals.
\end{project}

\begin{project}{Computational Geodynamo Physics-Informed Neural Network (PINN)}{Python, PyTorch, Optuna, PDEs}
  \item \textbf{Implemented} a PINN in PyTorch to solve a 2D magnetic induction equation (scalar advection-diffusion) as a toy model for the Earth's geodynamo.
  \item \textbf{Mathematically Derived} the total loss function by enforcing the PDE residual, initial Gaussian pulse conditions, and Dirichlet boundary conditions through automatic differentiation.
  \item \textbf{Optimised} the architecture and hyperparameters using an Optuna study to minimise the physics-informed residual loss across the spatial-temporal domain.
\end{project}

\begin{project}{Facial Recognition Toolkit}{Python, OpenCV, face\_recognition (dlib)}
  \item \textbf{Engineered} a Python command-line tool using OpenCV and the \texttt{face\_recognition}
        library to perform real-time face detection and identification via both static images and
        live webcam streams.
  \item \textbf{Implemented} 128-dimensional face-embedding extraction (HOG + CNN pipeline)
        with Euclidean-distance matching, supporting configurable tolerance thresholds for
        precision/recall trade-off in low-light conditions.
\end{project}

% ════════════════════════════════════════════════════════════
%  PROJECTS — NUMERICAL SIMULATIONS & COMPUTATIONAL PHYSICS
% ════════════════════════════════════════════════════════════
\section{Projects: Numerical Simulations \& Computational Physics}

\begin{project}{Quantum Mechanics: 1D Schrödinger Equation Solver}{Python, NumPy, SciPy (sparse eigensolvers), Matplotlib}
  \item \textbf{Implemented} a finite-difference numerical solver for the time-independent Schrödinger equation using sparse matrix eigenvalue methods.
  \item \textbf{Solved} multiple potential types including infinite/finite wells, harmonic oscillators (Hermite wavefunctions), and double wells.
  \item \textbf{Visualised} energy eigenvalues, wavefunction nodes, and quantum tunnelling phenomena with exponential transmission dependence on barrier width.
\end{project}

\begin{project}{Variational Quantum Monte Carlo (VQMC) Engine}{Rust, Metropolis Algorithm, Slater-Jastrow wavefunctions}
  \item \textbf{Implemented} a high-performance VMC engine in Rust to estimate the ground-state energy of small atoms (Helium, Lithium) in atomic units.
  \item \textbf{Simulated} electron configurations using the Metropolis algorithm with pluggable trial wavefunctions (Hydrogenic, Pad\'{e}, and Slater-Jastrow).
  \item \textbf{Optimised} for multi-electron antisymmetry using Slater determinants ($D_{\uparrow}$/$D_{\downarrow}$) and visualised energy convergence via CSV trace logging.
\end{project}

\begin{project}{High-Performance Fluid Dynamics Simulator (Lattice Boltzmann D2Q9)}{C++17, LBM, CMake}
  \item \textbf{Programmed} a D2Q9 Lattice Boltzmann Method simulation in C++ using the BGK collision model.
  \item \textbf{Implemented} the ``stream and collide'' algorithm with periodic boundary conditions to simulate fluid flow.
\end{project}

\begin{project}{Barium Cloud Diffusion in the Ionosphere}{Python, NumPy, SciPy, Monte Carlo sensitivity analysis}
  \item \textbf{Modelled} barium cloud diffusion in the lower ionosphere by numerically solving the
        diffusion PDE $\partial_t n = \nabla\cdot(D\nabla n)$ using both 1D and 3D explicit
        finite-difference (FTCS) schemes with Dirichlet and Neumann boundary conditions.
  \item \textbf{Implemented} semi-analytical integration using Bessel function expansions to provide
        a high-accuracy reference solution against which the numerical schemes were validated,
        confirming second-order spatial convergence.
  \item \textbf{Simulated} a Monte Carlo sensitivity analysis over the diffusion coefficient $D$
        and initial cloud parameters to extract atmospheric quantities from synthetic observation
        data, propagating uncertainties through the non-linear inversion pipeline.
\end{project}

\begin{project}{Z Boson Rediscovery --- CERN Open Data Analysis}{Python, Uproot, Awkward Array, SciPy, Matplotlib}
  \item \textbf{Processed} real proton-proton collision event files from the CERN Open Data Portal
        using \texttt{Uproot} and \texttt{Awkward Array} to reconstruct dimuon and dielectron
        invariant mass spectra from four-vectors:
        $M^2 = (E_1+E_2)^2 - |\mathbf{p}_1+\mathbf{p}_2|^2c^2$.
  \item \textbf{Rediscovered} the Z boson resonance at $M_Z \approx 91.2\,\mathrm{GeV}/c^2$ by
        fitting a Breit-Wigner (relativistic) line shape convolved with a Gaussian detector
        resolution function using \texttt{scipy.optimize.curve\_fit}.
  \item \textbf{Visualised} the full invariant mass spectrum from $20\,\mathrm{GeV}$ to
        $120\,\mathrm{GeV}$, identifying background contributions and comparing the measured
        decay width $\Gamma_Z$ against the PDG value.
\end{project}

\begin{project}{Ising Model Monte Carlo Simulation}{Python, NumPy, Metropolis algorithm, statistical mechanics}
  \item \textbf{Simulated} the Ising model of ferromagnetism on 2D lattices (square, triangular, hexagonal) using the Metropolis algorithm.
  \item \textbf{Visualised} real-time spin evolution and identified phase transitions at the critical temperature ($T_c \approx 2.269$ for square lattices) through thermodynamic analysis of magnetization and specific heat.
\end{project}

% ════════════════════════════════════════════════════════════
%  PROJECTS — QUANTITATIVE FINANCE & ALGORITHMIC TRADING
% ════════════════════════════════════════════════════════════
\section{Projects: Quantitative Finance \& Algorithmic Trading}

\begin{project}{Ombra Algorithmic Trading Framework}{Python, HMM, SQLite, Optuna, Streamlit, OANDA API}
  \item \textbf{Engineered} a multi-asset institutional-grade trading bot incorporating Hidden
        Markov Model (HMM) regime detection to classify market states (trending, mean-reverting,
        volatile), dynamically adjusting position sizing and entry logic per regime.
  \item \textbf{Implemented} limit-order execution with configurable slippage guardrails and
        automated SQLite database pruning to maintain a rolling history without unbounded
        storage growth, supporting continuous 24/7 operation.
  \item \textbf{Optimised} strategy hyperparameters (EMA spans, HMM states, risk fraction) via
        Optuna on historical replay backtests, maximising the out-of-sample Sharpe ratio and
        minimising maximum drawdown; deployed a Streamlit telemetry dashboard for live
        equity-curve monitoring.
\end{project}

\begin{project}{OANDA Trend Analysis Suite --- ``Trendy''}{Python, OANDA REST API, pandas, Matplotlib}
  \item \textbf{Engineered} a multi-asset trend screener analysing 44 instruments (Forex, Commodities, Indices) using the OANDA REST API.
  \item \textbf{Designed} a sophisticated EMA-based algorithm (8, 50, and 200 EMA) with multiple timeframe alignment (e.g., 4hr vs.\ weekly) to identify long/short trends.
  \item \textbf{Optimised} for speed by creating simplified versions (up to 75\% faster) and packaged the suite into standalone Windows 64-bit executables.
\end{project}

% ════════════════════════════════════════════════════════════
%  PROJECTS — SOFTWARE ENGINEERING
% ════════════════════════════════════════════════════════════
\section{Projects: Full-Stack \& Mobile Software Engineering}

\begin{project}{GameifyLife --- Hybrid Academic Dashboard}{Next.js, Capacitor, Gemini 2.0 API, TypeScript}
  \item \textbf{Engineered} a hybrid mobile app (iOS/Android) using Next.js and Capacitor, featuring a Gemini 2.0 Flash AI assistant for smart study recommendations.
  \item \textbf{Implemented} 15+ native features including haptic feedback, camera integration, GPS tracking, and background synchronisation.
\end{project}

\begin{project}{Turf Runner --- Location-Based Multiplayer Territory Game}{Flutter, Dart, Node.js, PostGIS, WebSockets}
  \item \textbf{Developed} a Flutter and Node.js real-time multiplayer mobile game in which players
        capture geo-fenced real-world territories, with game-state managed server-side via
        WebSocket connections for sub-100\,ms latency updates.
  \item \textbf{Engineered} conflict resolution and closed-loop territory validation using the
        Haversine formula for great-circle distance and PostGIS spatial queries
        (\texttt{ST\_Contains}, \texttt{ST\_Intersects}) on a PostgreSQL backend.
  \item \textbf{Simulated} concurrent multi-player sessions with automated Jest integration tests,
        verifying territory-ownership consistency under race-condition scenarios.
\end{project}

\begin{project}{MotionMonitor (Vesta) --- Clinical Fall Detection App}{React Native, Flask, Socket.IO, Python, CSV telemetry}
  \item \textbf{Built} a React Native mobile application streaming live 6-DoF IMU telemetry
        (accelerometer + gyroscope) at up to 100\,Hz to a Flask backend via Socket.IO,
        with sub-frame latency suitable for real-time clinical monitoring.
  \item \textbf{Implemented} a fall-detection alert loop using threshold-based peak detection on
        the resultant acceleration vector $|\mathbf{a}|=\sqrt{a_x^2+a_y^2+a_z^2}$,
        triggering push notifications and clinician-dashboard alerts on fall events.
  \item \textbf{Engineered} a CSV telemetry-logging pipeline to accumulate labelled sensor data
        for future supervised ML model development (e.g., LSTM-based fall classification).
\end{project}

\begin{project}{Real-Time Multiplayer Quiz Platform}{React (Vite), Node.js, Socket.io, TypeScript}
  \item \textbf{Engineered} a low-latency multiplayer quiz platform supporting up to 12 concurrent
        players per room via Socket.io namespaces, with server-authoritative game-state management
        to prevent client-side cheating.
  \item \textbf{Implemented} room lifecycle management (create, join, auto-expire), real-time
        leaderboard updates on each answer submission, and adaptive question timing based on
        per-room difficulty settings.
\end{project}

\begin{project}{Delivery Route Optimiser --- ``Path Finder''}{Python, OR-Tools, Google Maps API, OpenRouteService, Genetic Algorithms}
  \item \textbf{Implemented} route optimisation for multiple addresses using OSRM, Google Maps, and OpenRouteService with automatic fallback.
  \item \textbf{Engineered} solvers for the Travelling Salesman Problem (TSP) using Nearest Neighbour, 2-opt improvement, OR-Tools, and Genetic Algorithms.
\end{project}

% ════════════════════════════════════════════════════════════
%  ADDITIONAL ACADEMIC COURSEWORK & SKILLS
% ════════════════════════════════════════════════════════════
\section{Additional Coursework \& Academic Achievements}

\begin{itemize}
  \item \textbf{Advanced Quantum Mechanics:} Perturbation theory (time-independent and time-dependent),
        variational principle, WKB approximation, identical particles and Pauli exclusion.
  \item \textbf{General Relativity:} Tensor calculus, covariant differentiation, geodesic equation,
        Schwarzschild and Kerr metrics, Penrose diagrams, Hawking radiation (conceptual).
  \item \textbf{Statistical Mechanics:} Canonical and grand canonical ensembles, partition functions,
        phase transitions, critical phenomena, Bose-Einstein and Fermi-Dirac statistics.
  \item \textbf{Particle Physics:} Standard Model overview, Feynman diagram calculation (tree-level QED),
        hadron spectroscopy, detector physics, data analysis with ROOT/Uproot.
  \item \textbf{Electrodynamics:} Maxwell's equations in covariant form, Lienard-Wiechert potentials,
        radiation from accelerating charges, waveguides and cavities.
  \item \textbf{Mathematical Methods:} Complex analysis (contour integration, residue theorem),
        Fourier and Laplace transforms, Green's functions, group theory fundamentals.
\end{itemize}

% ════════════════════════════════════════════════════════════
%  LANGUAGES & SOFT SKILLS
% ════════════════════════════════════════════════════════════
\section{Languages \& Transferable Skills}

\begin{multicols}{2}
\begin{itemize}
  \item English (Native)
  \item Hindi / Punjabi (Fluent)
  \item Technical report writing \& \LaTeX\ typesetting
  \item Scientific communication and data storytelling
  \item Agile development (Git branching, PR workflows)
  \item Independent research and rapid upskilling
\end{itemize}
\end{multicols}

% ════════════════════════════════════════════════════════════
%  INTERESTS
% ════════════════════════════════════════════════════════════
\section{Interests}

Gravitational wave astronomy and next-generation detector design $\cdot$
Quantum computing algorithms for physics simulation $\cdot$
High-performance scientific computing (GPU acceleration, CUDA) $\cdot$
Competitive mathematics and physics olympiads $\cdot$
Open-source scientific software contribution

\vfill
{\small\color{light}\textit{References available on request. CV prepared with \LaTeX. Last updated: February 2026.}}

\end{document}
